% Options for packages loaded elsewhere
\PassOptionsToPackage{unicode}{hyperref}
\PassOptionsToPackage{hyphens}{url}
\PassOptionsToPackage{dvipsnames,svgnames,x11names}{xcolor}
\documentclass[
  11pt,
  a4paper,
]{article}
\usepackage{xcolor}
\usepackage[margin=1in]{geometry}
\usepackage{amsmath,amssymb}
\setcounter{secnumdepth}{-\maxdimen} % remove section numbering
\usepackage{iftex}
\ifPDFTeX
  \usepackage[T1]{fontenc}
  \usepackage[utf8]{inputenc}
  \usepackage{textcomp} % provide euro and other symbols
\else % if luatex or xetex
  \usepackage{unicode-math} % this also loads fontspec
  \defaultfontfeatures{Scale=MatchLowercase}
  \defaultfontfeatures[\rmfamily]{Ligatures=TeX,Scale=1}
\fi
\usepackage{lmodern}
\ifPDFTeX\else
  % xetex/luatex font selection
  \ifLuaTeX
    \usepackage{luaotfload}
    \directlua{luaotfload.add_fallback("mainfontfallback",{
      "NotoSerifBengali:mode=harf"
    })}
  \fi
  \setmainfont[,RawFeature={fallback=mainfontfallback}]{Latin Modern Roman}
  \ifLuaTeX
    \usepackage{luaotfload}
    \directlua{luaotfload.add_fallback("sansfontfallback",{
      "NotoSansBengali:mode=harf"
    })}
  \fi
  \setsansfont[,RawFeature={fallback=sansfontfallback}]{Latin Modern Sans}
  \ifLuaTeX
    \usepackage{luaotfload}
    \directlua{luaotfload.add_fallback("monofontfallback",{
      "NotoSansBengali:mode=harf"
    })}
  \fi
  \setmonofont[,RawFeature={fallback=monofontfallback}]{DejaVu Sans Mono}
\fi
% Use upquote if available, for straight quotes in verbatim environments
\IfFileExists{upquote.sty}{\usepackage{upquote}}{}
\IfFileExists{microtype.sty}{% use microtype if available
  \usepackage[]{microtype}
  \UseMicrotypeSet[protrusion]{basicmath} % disable protrusion for tt fonts
}{}
\makeatletter
\@ifundefined{KOMAClassName}{% if non-KOMA class
  \IfFileExists{parskip.sty}{%
    \usepackage{parskip}
  }{% else
    \setlength{\parindent}{0pt}
    \setlength{\parskip}{6pt plus 2pt minus 1pt}}
}{% if KOMA class
  \KOMAoptions{parskip=half}}
\makeatother
\usepackage{color}
\usepackage{fancyvrb}
\newcommand{\VerbBar}{|}
\newcommand{\VERB}{\Verb[commandchars=\\\{\}]}
\DefineVerbatimEnvironment{Highlighting}{Verbatim}{commandchars=\\\{\}}
% Add ',fontsize=\small' for more characters per line
\usepackage{framed}
\definecolor{shadecolor}{RGB}{248,248,248}
\newenvironment{Shaded}{\begin{snugshade}}{\end{snugshade}}
\newcommand{\AlertTok}[1]{\textcolor[rgb]{0.94,0.16,0.16}{#1}}
\newcommand{\AnnotationTok}[1]{\textcolor[rgb]{0.56,0.35,0.01}{\textbf{\textit{#1}}}}
\newcommand{\AttributeTok}[1]{\textcolor[rgb]{0.13,0.29,0.53}{#1}}
\newcommand{\BaseNTok}[1]{\textcolor[rgb]{0.00,0.00,0.81}{#1}}
\newcommand{\BuiltInTok}[1]{#1}
\newcommand{\CharTok}[1]{\textcolor[rgb]{0.31,0.60,0.02}{#1}}
\newcommand{\CommentTok}[1]{\textcolor[rgb]{0.56,0.35,0.01}{\textit{#1}}}
\newcommand{\CommentVarTok}[1]{\textcolor[rgb]{0.56,0.35,0.01}{\textbf{\textit{#1}}}}
\newcommand{\ConstantTok}[1]{\textcolor[rgb]{0.56,0.35,0.01}{#1}}
\newcommand{\ControlFlowTok}[1]{\textcolor[rgb]{0.13,0.29,0.53}{\textbf{#1}}}
\newcommand{\DataTypeTok}[1]{\textcolor[rgb]{0.13,0.29,0.53}{#1}}
\newcommand{\DecValTok}[1]{\textcolor[rgb]{0.00,0.00,0.81}{#1}}
\newcommand{\DocumentationTok}[1]{\textcolor[rgb]{0.56,0.35,0.01}{\textbf{\textit{#1}}}}
\newcommand{\ErrorTok}[1]{\textcolor[rgb]{0.64,0.00,0.00}{\textbf{#1}}}
\newcommand{\ExtensionTok}[1]{#1}
\newcommand{\FloatTok}[1]{\textcolor[rgb]{0.00,0.00,0.81}{#1}}
\newcommand{\FunctionTok}[1]{\textcolor[rgb]{0.13,0.29,0.53}{\textbf{#1}}}
\newcommand{\ImportTok}[1]{#1}
\newcommand{\InformationTok}[1]{\textcolor[rgb]{0.56,0.35,0.01}{\textbf{\textit{#1}}}}
\newcommand{\KeywordTok}[1]{\textcolor[rgb]{0.13,0.29,0.53}{\textbf{#1}}}
\newcommand{\NormalTok}[1]{#1}
\newcommand{\OperatorTok}[1]{\textcolor[rgb]{0.81,0.36,0.00}{\textbf{#1}}}
\newcommand{\OtherTok}[1]{\textcolor[rgb]{0.56,0.35,0.01}{#1}}
\newcommand{\PreprocessorTok}[1]{\textcolor[rgb]{0.56,0.35,0.01}{\textit{#1}}}
\newcommand{\RegionMarkerTok}[1]{#1}
\newcommand{\SpecialCharTok}[1]{\textcolor[rgb]{0.81,0.36,0.00}{\textbf{#1}}}
\newcommand{\SpecialStringTok}[1]{\textcolor[rgb]{0.31,0.60,0.02}{#1}}
\newcommand{\StringTok}[1]{\textcolor[rgb]{0.31,0.60,0.02}{#1}}
\newcommand{\VariableTok}[1]{\textcolor[rgb]{0.00,0.00,0.00}{#1}}
\newcommand{\VerbatimStringTok}[1]{\textcolor[rgb]{0.31,0.60,0.02}{#1}}
\newcommand{\WarningTok}[1]{\textcolor[rgb]{0.56,0.35,0.01}{\textbf{\textit{#1}}}}
\usepackage{longtable,booktabs,array}
\usepackage{calc} % for calculating minipage widths
% Correct order of tables after \paragraph or \subparagraph
\usepackage{etoolbox}
\makeatletter
\patchcmd\longtable{\par}{\if@noskipsec\mbox{}\fi\par}{}{}
\makeatother
% Allow footnotes in longtable head/foot
\IfFileExists{footnotehyper.sty}{\usepackage{footnotehyper}}{\usepackage{footnote}}
\makesavenoteenv{longtable}
\setlength{\emergencystretch}{3em} % prevent overfull lines
\providecommand{\tightlist}{%
  \setlength{\itemsep}{0pt}\setlength{\parskip}{0pt}}
% paper_preamble.tex — shared LaTeX preamble for the PCMA/SI paper set.
%
% Compiled with lualatex (NOT pdflatex) because several papers contain
% Bengali script: the CLL paper realises one cognitive state into
% English, Bengali and Banglish, so the Bengali output is EVIDENCE and
% cannot be transliterated away without destroying the point being made.
%
% Font selection is NOT done here. pandoc >= 3 has native
% `mainfontfallback` / `monofontfallback` variables that emit the
% luaotfload fallback registration in the correct position relative to
% unicode-math; doing it by hand in this header lands after unicode-math
% and breaks font loading outright. See tools/build_papers.sh.

\usepackage{amsmath}
\usepackage{amssymb}
\usepackage{mathtools}   % \xleftrightarrow etc. used in cll.md §5
\usepackage{booktabs}

% ---------------------------------------------------------------------
% Unicode symbols appearing in PROSE rather than inside $...$.
%
% The papers were written as markdown and read as markdown, so they use
% real characters — "the ω-boost", "α-child", "ΔR coherence" — instead of
% LaTeX macros. Under unicode-math those codepoints are math-mode-only
% and abort the build ("\symbballowed only in math mode"). Mapping each
% to \ensuremath keeps the markdown readable on GitHub AND compilable,
% which is the whole point of keeping one source.
% ---------------------------------------------------------------------
\usepackage{newunicodechar}

% Greek, lower case
\newunicodechar{α}{\ensuremath{\alpha}}
\newunicodechar{β}{\ensuremath{\beta}}
\newunicodechar{γ}{\ensuremath{\gamma}}
\newunicodechar{δ}{\ensuremath{\delta}}
\newunicodechar{ε}{\ensuremath{\varepsilon}}
\newunicodechar{η}{\ensuremath{\eta}}
\newunicodechar{θ}{\ensuremath{\theta}}
\newunicodechar{κ}{\ensuremath{\kappa}}
\newunicodechar{λ}{\ensuremath{\lambda}}
\newunicodechar{μ}{\ensuremath{\mu}}
\newunicodechar{π}{\ensuremath{\pi}}
\newunicodechar{ρ}{\ensuremath{\rho}}
\newunicodechar{σ}{\ensuremath{\sigma}}
\newunicodechar{τ}{\ensuremath{\tau}}
\newunicodechar{φ}{\ensuremath{\varphi}}
\newunicodechar{χ}{\ensuremath{\chi}}
\newunicodechar{ψ}{\ensuremath{\psi}}
\newunicodechar{ω}{\ensuremath{\omega}}
\newunicodechar{ν}{\ensuremath{\nu}}
\newunicodechar{ξ}{\ensuremath{\xi}}
\newunicodechar{ζ}{\ensuremath{\zeta}}
\newunicodechar{ι}{\ensuremath{\iota}}

% Greek, upper case
\newunicodechar{Δ}{\ensuremath{\Delta}}
\newunicodechar{Φ}{\ensuremath{\Phi}}
\newunicodechar{Σ}{\ensuremath{\Sigma}}
\newunicodechar{Ψ}{\ensuremath{\Psi}}
\newunicodechar{Ω}{\ensuremath{\Omega}}
\newunicodechar{Γ}{\ensuremath{\Gamma}}
\newunicodechar{Λ}{\ensuremath{\Lambda}}
\newunicodechar{Θ}{\ensuremath{\Theta}}
\newunicodechar{Π}{\ensuremath{\Pi}}

% Operators and relations
\newunicodechar{×}{\ensuremath{\times}}
\newunicodechar{·}{\ensuremath{\cdot}}
\newunicodechar{±}{\ensuremath{\pm}}
\newunicodechar{−}{\ensuremath{-}}
\newunicodechar{≈}{\ensuremath{\approx}}
\newunicodechar{≤}{\ensuremath{\leq}}
\newunicodechar{≥}{\ensuremath{\geq}}
\newunicodechar{≠}{\ensuremath{\neq}}
\newunicodechar{∈}{\ensuremath{\in}}
\newunicodechar{∂}{\ensuremath{\partial}}
\newunicodechar{∇}{\ensuremath{\nabla}}
\newunicodechar{√}{\ensuremath{\sqrt{}}}
\newunicodechar{∞}{\ensuremath{\infty}}
\newunicodechar{∝}{\ensuremath{\propto}}
\newunicodechar{∑}{\ensuremath{\sum}}
\newunicodechar{∫}{\ensuremath{\int}}
\newunicodechar{ℝ}{\ensuremath{\mathbb{R}}}
\newunicodechar{ℕ}{\ensuremath{\mathbb{N}}}
\newunicodechar{‖}{\ensuremath{\|}}
\newunicodechar{⊕}{\ensuremath{\oplus}}
\newunicodechar{⊗}{\ensuremath{\otimes}}
\newunicodechar{∀}{\ensuremath{\forall}}
\newunicodechar{∃}{\ensuremath{\exists}}
\newunicodechar{¬}{\ensuremath{\neg}}
\newunicodechar{∧}{\ensuremath{\wedge}}
\newunicodechar{∨}{\ensuremath{\vee}}

% Arrows
\newunicodechar{→}{\ensuremath{\rightarrow}}
\newunicodechar{←}{\ensuremath{\leftarrow}}
\newunicodechar{↔}{\ensuremath{\leftrightarrow}}
\newunicodechar{↑}{\ensuremath{\uparrow}}
\newunicodechar{↓}{\ensuremath{\downarrow}}
\newunicodechar{⇒}{\ensuremath{\Rightarrow}}
\newunicodechar{⇐}{\ensuremath{\Leftarrow}}
\newunicodechar{⇔}{\ensuremath{\Leftrightarrow}}

% Super/subscripts written as literal characters
\newunicodechar{²}{\ensuremath{^{2}}}
\newunicodechar{³}{\ensuremath{^{3}}}
\newunicodechar{⁴}{\ensuremath{^{4}}}
\newunicodechar{⁵}{\ensuremath{^{5}}}
\newunicodechar{⁶}{\ensuremath{^{6}}}
\newunicodechar{₀}{\ensuremath{_{0}}}
\newunicodechar{₁}{\ensuremath{_{1}}}
\newunicodechar{₂}{\ensuremath{_{2}}}
\newunicodechar{₃}{\ensuremath{_{3}}}
\newunicodechar{ᵀ}{\ensuremath{^{\mathsf{T}}}}

% Status glyphs. The papers use emoji as status markers in tables
% (implemented / failed / warning). Mapped to typographic equivalents:
% a colour emoji font in lualatex is unreliable, and the MEANING is what
% carries — a reader needs to see "this one failed", not a red cross.
\newunicodechar{✅}{\ensuremath{\checkmark}}
\newunicodechar{✓}{\ensuremath{\checkmark}}
\newunicodechar{❌}{\ensuremath{\times}}
\newunicodechar{✗}{\ensuremath{\times}}
\newunicodechar{⚠}{\textbf{!}}
\newunicodechar{🧩}{\textbf{[piece]}}
\newunicodechar{📋}{\textbf{[list]}}
\newunicodechar{🟡}{\textbf{[partial]}}
\newunicodechar{🔴}{\textbf{[dead]}}
\newunicodechar{🟢}{\textbf{[live]}}
\newunicodechar{▼}{\ensuremath{\blacktriangledown}}
\newunicodechar{►}{\ensuremath{\blacktriangleright}}
\newunicodechar{□}{\ensuremath{\square}}
\newunicodechar{█}{\rule{0.6em}{0.9ex}}
\newunicodechar{•}{\textbullet}
\newunicodechar{†}{\textdagger}
\newunicodechar{‑}{-}
\newunicodechar{‐}{-}
\newunicodechar{…}{\ldots}
\newunicodechar{।}{\textbar}   % danda: rendered by the Bengali fallback in text
\newunicodechar{⁻}{\ensuremath{^{-}}}

% Zero-width formatting codepoints. These are invisible in an editor but
% are real characters to LaTeX: an emoji written as "⚠️" is TWO
% codepoints (U+26A0 U+FE0F), so declaring only the visible one leaves a
% variation selector with no definition and the build dies on a symbol
% nobody can see. Both are mapped to nothing.
\newunicodechar{^^^^fe0f}{}    % VARIATION SELECTOR-16
\newunicodechar{^^^^200d}{}    % ZERO WIDTH JOINER
\newunicodechar{^^^^200b}{}    % ZERO WIDTH SPACE

\usepackage{longtable}
\usepackage{array}

% Comparison matrices and cost ledgers are wide. Let them shrink to the
% text block rather than overflow the margin.
\usepackage{adjustbox}

% Code blocks: wrap long lines instead of running off the page. Several
% papers quote real C++ with long identifiers.
\usepackage{fvextra}
\fvset{breaklines=true,breakanywhere=true,fontsize=\small}

\setlength{\parskip}{0.5em}
\setlength{\parindent}{0pt}

\usepackage{fancyhdr}
\pagestyle{fancy}
\fancyhf{}
\fancyfoot[C]{\thepage}
\fancyhead[L]{\small\itshape Synthetic Intelligence — research lineage}
\fancyhead[R]{\small\itshape irfan@furylogic.com}
\renewcommand{\headrulewidth}{0.2pt}
\setlength{\headheight}{14pt}
\setlength{\headsep}{18pt}
\usepackage{bookmark}
\IfFileExists{xurl.sty}{\usepackage{xurl}}{} % add URL line breaks if available
\urlstyle{same}
\hypersetup{
  pdftitle={RB+G: Gated Recursive Bifurcation with Analogical Bridging --- A Physics Substrate for Contradiction Metabolism and Cross-Domain Repurposing},
  pdfauthor={Irfan Mahir --- Independent Researcher, Dhaka, Bangladesh --- irfan@furylogic.com},
  colorlinks=true,
  linkcolor={black},
  filecolor={Maroon},
  citecolor={Blue},
  urlcolor={[HTML]{0B5394}},
  pdfcreator={LaTeX via pandoc}}

\title{RB+G: Gated Recursive Bifurcation with Analogical Bridging --- A
Physics Substrate for Contradiction Metabolism and Cross-Domain
Repurposing}
\author{Irfan Mahir --- Independent Researcher, Dhaka, Bangladesh ---
irfan@furylogic.com}
\date{}

\begin{document}
\maketitle

{
\hypersetup{linkcolor=}
\setcounter{tocdepth}{3}
\tableofcontents
}
\section{RB+G: Gated Recursive Bifurcation with Analogical Bridging ---
A Physics Substrate for Contradiction Metabolism and Cross-Domain
Repurposing}\label{rbg-gated-recursive-bifurcation-with-analogical-bridging-a-physics-substrate-for-contradiction-metabolism-and-cross-domain-repurposing}

\textbf{Irfan Mahir} \emph{Independent Researcher, Dhaka, Bangladesh}
irfan@furylogic.com

\begin{quote}
\emph{``A concept that contains a contradiction is not a concept --- it
is two concepts waiting for the physics to split them.''}
\end{quote}

\begin{center}\rule{0.5\linewidth}{0.5pt}\end{center}

\subsection{Abstract}\label{abstract}

We present \textbf{RB+G} (Gated Recursive Bifurcation), a
physics-substrate extension of Resonant Bifurcation {[}Mahir 2026o{]}
that addresses four documented failure modes of the original mechanism
--- cascade fission, timeout under stress-test streams, null child
labels, and drive-amplitude confounds --- while retaining the core
insight: when a concept-cavity's positive and negative Hebbian updates
destructively interfere, the physics engine physically splits the cavity
into two children rather than balancing weights on a single site. The
extension introduces five mechanisms: (i) a generation cap that prevents
cascade, (ii) a parallel alias ledger that avoids graph-label pollution,
(iii) an energy-conserving \(1/\sqrt{2}\) probe fan-out that preserves
total injected energy across fissioned parents, (iv) a learned
\(\alpha \leftrightarrow \beta\) binding tunnel that replaces the fiat
coupling of the original, and (v) an explicit audit trail. On top of
RB+G we build \textbf{L6-A} (L6 Analogical Bridge), a discovery
mechanism that inserts a low-stiffness ``rubber'' tunnel between the
excitatory α-children of two fissioned concepts whenever they share
\(\geq K\) downstream excitatory targets. In the framework of
Constructive Structural Cognition {[}Mahir 2026p{]}, RB+G realizes Axiom
II (non-equilibrium phase-space expansion --- the graph grows to house
contradictions) and L6-A realizes Axiom III (active construction ---
analogies are physics-inserted, not retrieved). We validate the combined
pipeline with a 20-seed reproducibility study on six concepts (three
contradiction-cluster targets, one mono-directional control, two null
controls), a four-lever ablation table, and a drug-repurposing case
study on the insulin-signalling ontology in which a 3 Hz probe on
\texttt{metformin} surfaces \texttt{atorvastatin} --- a concept
unreachable by naive graph propagation but reachable through the
physics-inserted
\texttt{metformin.α\ ↔\ berberine.α\ ↔\ hmg\_coa\_reductase\ ↔\ atorvastatin}
bridge.

\begin{center}\rule{0.5\linewidth}{0.5pt}\end{center}

\subsection{1. Introduction}\label{introduction}

\subsubsection{1.1 The Contradiction Problem in Physics-Based
Cognition}\label{the-contradiction-problem-in-physics-based-cognition}

The PRM + MIMS substrate {[}Mahir 2026q, 2026r{]} represents concepts as
mass-spring cavities and relationships as physical tunnels with
material-dependent stiffness. When two tunnels of opposite polarity
(\(p_1 = +1\), \(p_2 = -1\)) drive the same cavity, their forces sum.
Over consolidation this produces destructive interference in the
ZeroGrad Hebbian update rule: positive and negative \(\Delta_e\) terms
cancel, and the effective plasticity approaches zero even though the
concept is under high tension. The cavity's ability to represent the
concept degrades.

This is the failure mode identified in \texttt{srg.md} and
\texttt{rb.md}: a concept like \texttt{metformin} has one incident
excitatory tunnel (\texttt{activates\ ampk}) and one incident inhibitory
tunnel (\texttt{inhibits\ mtor}). Under training, positive and negative
Hebbian updates land on the same cavity and cancel. The concept becomes
structurally invisible.

Two solutions have been proposed:

\begin{itemize}
\tightlist
\item
  \textbf{SRG} (Synaptic Receptor Gating, \texttt{srg.md}): keep the
  cavity as one entity but route excitatory and inhibitory forces into
  separate E-port and I-port storage. Waves entering the I-port are 180°
  phase-inverted before hitting the mass; the observable displacement is
  \(x = x_E - x_I\). Empirically SRG produced marginal regression on the
  metformin-recall benchmark (0.887 → 0.831); we retain it as an
  optional per-tick mechanism but do not depend on it here.
\item
  \textbf{RB} (Resonant Bifurcation, \texttt{rb.md}): physically split
  the cavity into two children --- metformin.α inheriting the excitatory
  tunnels, metformin.β inheriting the inhibitory tunnels --- bound by an
  ultra-stiff (k=2.5) semantic spring. This is the foundation on which
  RB+G is built.
\end{itemize}

\subsubsection{1.2 The Four Documented RB Failure
Modes}\label{the-four-documented-rb-failure-modes}

An honest implementation of the RB proposal on our substrate exhibited
four failure modes:

\begin{itemize}
\tightlist
\item
  \textbf{F1 Cascade.} With trigger threshold \(\tau > 0.005\) and no
  generation cap, fission propagated. On a linux-null control the engine
  fired 13 unintended fissions in one training run; on the luna 2022
  stress-test it timed out mid-stream.
\item
  \textbf{F2 Timeout.} The luna stress-test (18 daily rows from
  \texttt{data/luna\_2022\_may.csv}) triggered by the same cascade
  behaviour, exceeding the test harness's time budget.
\item
  \textbf{F3 Null child labels.} New cavities created by fission did not
  exist in \texttt{si.graph.by\_label} --- the label-to-cavity map was
  frozen at bind time --- so downstream readouts printed
  \texttt{\textless{}null\textgreater{}} for α and β children.
\item
  \textbf{F4 Drive-amplitude confound.} After fission, a drive on the
  parent hit both children at full amplitude, silently injecting
  \(2\times\) the energy of a single-cavity drive. Recall gains were
  confounded with the extra injected energy.
\end{itemize}

RB+G fixes all four while retaining the core mechanism.

\subsubsection{1.3 Contributions}\label{contributions}

\begin{itemize}
\tightlist
\item
  \textbf{Five-mechanism RB+G architecture} (§3) --- the four fixes plus
  an audit trail.
\item
  \textbf{Layer-2 fission trigger} (§4) --- \(I > 0.15\),
  \(\tau > 0.005\), sustained 3 ticks, generation \(< g_\text{max}\), at
  most one fire per structural tick.
\item
  \textbf{L6-A extension} (§5) --- physics-inserted rubber tunnels
  between α-siblings sharing downstream excitatory targets.
\item
  \textbf{20-seed reproducibility on six concepts} (§6.1) with three
  contradiction targets, one mono-directional control, two null
  controls.
\item
  \textbf{Four-lever ablation table} (§6.2) --- honest ablation, with
  two configurations producing \emph{higher} raw RMS than the full
  engine, both documented as measurement confounds.
\item
  \textbf{Drug-repurposing case study} (§6.3) --- a 3 Hz probe on
  \texttt{metformin} surfaces \texttt{atorvastatin} at rank 10 only
  after L6-A bridge insertion; the concept is unreachable by naive graph
  propagation.
\end{itemize}

\begin{center}\rule{0.5\linewidth}{0.5pt}\end{center}

\subsection{2. Position in the Prior
Lineage}\label{position-in-the-prior-lineage}

RB+G is a Layer-1 mechanism in the tri-layer architecture of
\texttt{synthetic-intelligence.md}. It does not replace anything; it
\emph{adds} structural mutability to the physics substrate.

\begin{longtable}[]{@{}
  >{\raggedright\arraybackslash}p{(\linewidth - 4\tabcolsep) * \real{0.3333}}
  >{\raggedright\arraybackslash}p{(\linewidth - 4\tabcolsep) * \real{0.3333}}
  >{\raggedright\arraybackslash}p{(\linewidth - 4\tabcolsep) * \real{0.3333}}@{}}
\toprule\noalign{}
\begin{minipage}[b]{\linewidth}\raggedright
Component
\end{minipage} & \begin{minipage}[b]{\linewidth}\raggedright
Provides
\end{minipage} & \begin{minipage}[b]{\linewidth}\raggedright
RB+G uses it for
\end{minipage} \\
\midrule\noalign{}
\endhead
\bottomrule\noalign{}
\endlastfoot
PRM (mass-spring cavities) {[}Mahir 2026q{]} & The cavity/tunnel
substrate & The medium into which children are spawned \\
MIMS (material impedance) {[}Mahir 2026r{]} & Per-relation stiffness /
damping & Sets the excitatory-vs-inhibitory tunnel polarity that
determines child routing \\
HNCA (Kuramoto macro-nodes) {[}Mahir 2026n{]} & Regional coupling for
the substrate & Children inherit their parent's cluster; macro-node
membership is extended \\
NTDA (D/C/S modulators) {[}Mahir 2026s{]} & Neurotransmitter scalars &
Training NT profile (\texttt{S=2}) accelerates Hebbian → drives
incoherence → triggers fission \\
ZeroGrad Hebbian {[}Mahir 2026d{]} & Local plasticity & The Hebbian
\(\Delta_e\) split into \(\Sigma^+\) and \(\Sigma^-\) per cavity is the
fission trigger signal \\
RIFA IC + SIL {[}Mahir 2026g{]} & Contradiction metabolism & Global
tension \(T(t)\) gates the fission budget (Axiom II) \\
Nine reasoning primitives {[}Mahir 2026p, §4.4{]} & Layer-2 composition
& The structural tick that computes \(I\) and \(\tau\) per cavity \\
Original RB {[}Mahir 2026o{]} & Fission concept & The base mechanism;
RB+G fixes its four failure modes \\
\end{longtable}

RB+G's novelty is not the \emph{idea} of fission (RB has that) but its
\emph{disciplined implementation} --- the five mechanisms of §3 that
make fission energy-conserving, cascade-safe, label-clean, and
auditable.

\begin{center}\rule{0.5\linewidth}{0.5pt}\end{center}

\subsection{3. The RB+G Architecture}\label{the-rbg-architecture}

\subsubsection{3.1 State Additions}\label{state-additions}

A cavity gains one integer:

\begin{Shaded}
\begin{Highlighting}[]
\KeywordTok{struct}\NormalTok{ Cavity }\OperatorTok{\{}
    \CommentTok{// ... existing PRM fields ...}
    \DataTypeTok{int}\NormalTok{   generation    }\OperatorTok{=} \DecValTok{0}\OperatorTok{;}   \CommentTok{// 0 for gen{-}0 (original graph cavities)}
\OperatorTok{\};}
\end{Highlighting}
\end{Shaded}

The engine gains four structures:

\begin{Shaded}
\begin{Highlighting}[]
\KeywordTok{enum} \KeywordTok{class}\NormalTok{ AliasRole }\OperatorTok{:} \DataTypeTok{int} \OperatorTok{\{}\NormalTok{ RootAlias }\OperatorTok{=} \DecValTok{0}\OperatorTok{,}\NormalTok{ Alpha }\OperatorTok{=} \DecValTok{1}\OperatorTok{,}\NormalTok{ Beta }\OperatorTok{=} \DecValTok{2} \OperatorTok{\};}
\KeywordTok{struct}\NormalTok{ AliasEntry }\OperatorTok{\{}
\NormalTok{    AliasRole role      }\OperatorTok{=}\NormalTok{ AliasRole}\OperatorTok{::}\NormalTok{RootAlias}\OperatorTok{;}
    \DataTypeTok{int}\NormalTok{       parent    }\OperatorTok{=} \OperatorTok{{-}}\DecValTok{1}\OperatorTok{;}   \CommentTok{// parent cavity id (self for RootAlias)}
    \DataTypeTok{int}\NormalTok{       partner   }\OperatorTok{=} \OperatorTok{{-}}\DecValTok{1}\OperatorTok{;}   \CommentTok{// sibling id (β if Alpha, α if Beta)}
    \DataTypeTok{int}\NormalTok{       label\_key }\OperatorTok{=} \OperatorTok{{-}}\DecValTok{1}\OperatorTok{;}   \CommentTok{// parent\textquotesingle{}s graph label index}
\OperatorTok{\};}
\KeywordTok{struct}\NormalTok{ FissionRecord }\OperatorTok{\{}
    \DataTypeTok{int}\NormalTok{    tick\_index}\OperatorTok{,}\NormalTok{ parent\_id}\OperatorTok{,}\NormalTok{ streak}\OperatorTok{,}\NormalTok{ child\_alpha}\OperatorTok{,}\NormalTok{ child\_beta}\OperatorTok{;}
    \DataTypeTok{int}\NormalTok{    incident\_alpha}\OperatorTok{,}\NormalTok{ incident\_beta}\OperatorTok{;}
    \DataTypeTok{float}\NormalTok{  I}\OperatorTok{,}\NormalTok{ tau}\OperatorTok{,}\NormalTok{ drive\_split\_factor}\OperatorTok{;}
    \BuiltInTok{std::}\NormalTok{string parent\_label}\OperatorTok{;}
\OperatorTok{\};}

\BuiltInTok{std::}\NormalTok{vector}\OperatorTok{\textless{}}\NormalTok{AliasEntry}\OperatorTok{\textgreater{}}\NormalTok{                       alias\_ledger}\OperatorTok{;}
\BuiltInTok{std::}\NormalTok{unordered\_map}\OperatorTok{\textless{}}\DataTypeTok{int}\OperatorTok{,}\BuiltInTok{std::}\NormalTok{pair}\OperatorTok{\textless{}}\DataTypeTok{int}\OperatorTok{,}\DataTypeTok{int}\OperatorTok{\textgreater{}\textgreater{}}\NormalTok{    fission\_map}\OperatorTok{;}
\BuiltInTok{std::}\NormalTok{vector}\OperatorTok{\textless{}}\NormalTok{FissionRecord}\OperatorTok{\textgreater{}}\NormalTok{                    fission\_records}\OperatorTok{;}
\DataTypeTok{int}   \VariableTok{g\_max}              \OperatorTok{=} \DecValTok{1}\OperatorTok{;}
\DataTypeTok{float}\NormalTok{ rbg\_bind\_k\_init    }\OperatorTok{=} \FloatTok{1.0}\BuiltInTok{f}\OperatorTok{;}
\DataTypeTok{float}\NormalTok{ rbg\_bind\_pol\_init  }\OperatorTok{=} \FloatTok{0.0}\BuiltInTok{f}\OperatorTok{;}
\DataTypeTok{float}\NormalTok{ rbg\_fanout\_factor  }\OperatorTok{=} \FloatTok{0.70710678}\BuiltInTok{f}\OperatorTok{;}   \CommentTok{// 1/√2}
\end{Highlighting}
\end{Shaded}

\subsubsection{3.2 The Five Mechanisms}\label{the-five-mechanisms}

\textbf{M1. Generation cap (fixes F1, F2).} Every child cavity is born
at \texttt{generation\ =\ parent.generation\ +\ 1}. The fission trigger
refuses any cavity with \texttt{generation\ \textgreater{}=\ g\_max}.
Setting \(g_\text{max} = 1\) means only original graph cavities can
fission, and their children cannot. This is a computational necessity,
not an accuracy win: cascade fission scatters energy across many small
cavities without improving concept-scoped recall (see §6.2).

\textbf{M2. Alias ledger (fixes F3).} The label-to-cavity map
\texttt{si.graph.by\_label} is frozen at bind. Instead of extending it,
RB+G maintains a parallel \texttt{alias\_ledger} sized to
\texttt{cavities.size()}. Every gen-0 cavity has a \texttt{RootAlias}
entry pointing at itself. On fission, the parent's ledger entry becomes
a \texttt{RootAlias} pointing at itself (with \texttt{label\_key}
retained), and two new entries are appended for α and β with
\texttt{label\_key\ =\ parent.label\_key}. Dumps that print synthetic
labels traverse the ledger:
\texttt{label\_of(α)\ =\ graph.nodes{[}label\_key{]}.label\ +\ ".α"}.

\textbf{M3. Energy-conserving fan-out (fixes F4).} The
\texttt{drive(node\_id,\ freq,\ amp,\ dur)} API checks the
\texttt{fission\_map}. If the driven cavity has been fissioned, the
drive is routed to both children at \(\text{amp} / \sqrt{2}\):

\begin{Shaded}
\begin{Highlighting}[]
\KeywordTok{auto}\NormalTok{ it }\OperatorTok{=}\NormalTok{ fission\_map}\OperatorTok{.}\NormalTok{find}\OperatorTok{(}\NormalTok{node\_id}\OperatorTok{);}
\ControlFlowTok{if} \OperatorTok{(}\NormalTok{it }\OperatorTok{!=}\NormalTok{ fission\_map}\OperatorTok{.}\NormalTok{end}\OperatorTok{()} \OperatorTok{\&\&}\NormalTok{ rbg\_alias\_route\_enabled}\OperatorTok{)} \OperatorTok{\{}
    \DataTypeTok{float}\NormalTok{ half\_amp }\OperatorTok{=}\NormalTok{ amp }\OperatorTok{*}\NormalTok{ rbg\_fanout\_factor}\OperatorTok{;}   \CommentTok{// 1/√2}
    \ControlFlowTok{for} \OperatorTok{(}\DataTypeTok{int}\NormalTok{ child }\OperatorTok{:} \OperatorTok{\{}\NormalTok{it}\OperatorTok{{-}\textgreater{}}\NormalTok{second}\OperatorTok{.}\NormalTok{first}\OperatorTok{,}\NormalTok{ it}\OperatorTok{{-}\textgreater{}}\NormalTok{second}\OperatorTok{.}\NormalTok{second}\OperatorTok{\})} \OperatorTok{\{}
\NormalTok{        cavities}\OperatorTok{[}\NormalTok{child}\OperatorTok{].}\NormalTok{drive\_freq     }\OperatorTok{=}\NormalTok{ freq}\OperatorTok{;}
\NormalTok{        cavities}\OperatorTok{[}\NormalTok{child}\OperatorTok{].}\NormalTok{drive\_amp      }\OperatorTok{=}\NormalTok{ half\_amp}\OperatorTok{;}
\NormalTok{        cavities}\OperatorTok{[}\NormalTok{child}\OperatorTok{].}\NormalTok{drive\_envelope }\OperatorTok{=} \FloatTok{1.0}\BuiltInTok{f}\OperatorTok{;}
\NormalTok{        cavities}\OperatorTok{[}\NormalTok{child}\OperatorTok{].}\NormalTok{drive\_decay    }\OperatorTok{=} \BuiltInTok{std::}\NormalTok{exp}\OperatorTok{({-}}\FloatTok{3.0}\BuiltInTok{f} \OperatorTok{/}\NormalTok{ d}\OperatorTok{);}
\NormalTok{        cavities}\OperatorTok{[}\NormalTok{child}\OperatorTok{].}\NormalTok{was\_driven     }\OperatorTok{=} \KeywordTok{true}\OperatorTok{;}
    \OperatorTok{\}}
    \ControlFlowTok{return}\OperatorTok{;}
\OperatorTok{\}}
\end{Highlighting}
\end{Shaded}

The parent cavity itself is zeroed at fission and becomes a ``routing
ghost'' --- it holds no field but its id remains the addressing target
so external callers do not need to know that fission happened. Total
kinetic energy injected by a probe on a fissioned parent equals a
single-cavity probe:
\(2 \cdot \frac{1}{2} m (\text{amp}/\sqrt{2})^2 = \frac{1}{2} m \, \text{amp}^2\).

\textbf{M4. Learned binding tunnel (replaces RB's fiat \(k=2.5\)).}
After the split, a new tunnel is inserted between α and β:

\begin{Shaded}
\begin{Highlighting}[]
\NormalTok{Tunnel bind}\OperatorTok{\{\};}
\NormalTok{bind}\OperatorTok{.}\NormalTok{cavity\_a       }\OperatorTok{=}\NormalTok{ alpha}\OperatorTok{;}
\NormalTok{bind}\OperatorTok{.}\NormalTok{cavity\_b       }\OperatorTok{=}\NormalTok{ beta}\OperatorTok{;}
\NormalTok{bind}\OperatorTok{.}\NormalTok{stiffness      }\OperatorTok{=}\NormalTok{ rbg\_bind\_k\_init}\OperatorTok{;}   \CommentTok{// 1.0}
\NormalTok{bind}\OperatorTok{.}\NormalTok{base\_stiffness }\OperatorTok{=}\NormalTok{ rbg\_bind\_k\_init}\OperatorTok{;}
\NormalTok{bind}\OperatorTok{.}\NormalTok{damping        }\OperatorTok{=} \FloatTok{0.020}\BuiltInTok{f}\OperatorTok{;}
\NormalTok{bind}\OperatorTok{.}\NormalTok{polarity       }\OperatorTok{=}\NormalTok{ rbg\_bind\_pol\_init}\OperatorTok{;}  \CommentTok{// 0.0 — pure coordination}
\NormalTok{bind}\OperatorTok{.}\NormalTok{plasticity\_delta }\OperatorTok{=} \FloatTok{0.0}\BuiltInTok{f}\OperatorTok{;}
\NormalTok{tunnels}\OperatorTok{.}\NormalTok{push\_back}\OperatorTok{(}\NormalTok{bind}\OperatorTok{);}
\end{Highlighting}
\end{Shaded}

Init \(k = 1.0\) (soft), polarity \(p = 0.0\) (no forced sign). Under
Hebbian consolidation the stiffness moves in \([0.5, 2.0]\); the
polarity shifts under P5b splitting toward whichever sign the
correlation actually prefers. Empirically converges to a coordination
stiffness that is \emph{lower} than RB's fiat 2.5 and equally effective
(see §6.2 ABLATE\_BIND).

\textbf{M5. Audit trail.} Every fission appends a \texttt{FissionRecord}
to \texttt{physics.fission\_records}. The \texttt{fission\_ledger} shell
command writes the vector to \texttt{data/fission\_ledger.jsonl} with
tick, parent label, \(I\), \(\tau\), streak, child ids, incident tunnel
counts, and drive split factor. Reproducibility is possible from the
ledger alone: given the axiom files, the RNG seed formula, and the
ledger, any run can be replayed.

\subsubsection{3.3 The Fission Operation}\label{the-fission-operation}

\begin{Shaded}
\begin{Highlighting}[]
\BuiltInTok{std::}\NormalTok{pair}\OperatorTok{\textless{}}\DataTypeTok{int}\OperatorTok{,}\DataTypeTok{int}\OperatorTok{\textgreater{}}\NormalTok{ fission\_locked}\OperatorTok{(}\DataTypeTok{int}\NormalTok{ parent}\OperatorTok{,} \DataTypeTok{float}\NormalTok{ I\_val}\OperatorTok{,} \DataTypeTok{float}\NormalTok{ tau\_val}\OperatorTok{,}
                                   \DataTypeTok{int}\NormalTok{ streak}\OperatorTok{,} \DataTypeTok{int}\NormalTok{ tick\_index}\OperatorTok{)} \OperatorTok{\{}
    \ControlFlowTok{if} \OperatorTok{(}\NormalTok{parent }\OperatorTok{\textless{}} \DecValTok{0} \OperatorTok{||} \OperatorTok{(}\DataTypeTok{size\_t}\OperatorTok{)}\NormalTok{parent }\OperatorTok{\textgreater{}=}\NormalTok{ cavities}\OperatorTok{.}\NormalTok{size}\OperatorTok{())} \ControlFlowTok{return} \OperatorTok{\{{-}}\DecValTok{1}\OperatorTok{,{-}}\DecValTok{1}\OperatorTok{\};}
    \ControlFlowTok{if} \OperatorTok{(}\NormalTok{fission\_map}\OperatorTok{.}\NormalTok{count}\OperatorTok{(}\NormalTok{parent}\OperatorTok{))} \ControlFlowTok{return}\NormalTok{ fission\_map}\OperatorTok{[}\NormalTok{parent}\OperatorTok{];}    \CommentTok{// idempotent}
    \ControlFlowTok{if} \OperatorTok{(}\NormalTok{cavities}\OperatorTok{[}\NormalTok{parent}\OperatorTok{].}\NormalTok{generation }\OperatorTok{\textgreater{}=} \VariableTok{g\_max}\OperatorTok{)} \ControlFlowTok{return} \OperatorTok{\{{-}}\DecValTok{1}\OperatorTok{,{-}}\DecValTok{1}\OperatorTok{\};}     \CommentTok{// M1}

\NormalTok{    Cavity p\_snap }\OperatorTok{=}\NormalTok{ cavities}\OperatorTok{[}\NormalTok{parent}\OperatorTok{];}
    \CommentTok{// zero the parent (routing ghost)}
\NormalTok{    cavities}\OperatorTok{[}\NormalTok{parent}\OperatorTok{]} \OperatorTok{=}\NormalTok{ zeroed}\OperatorTok{(}\NormalTok{p\_snap}\OperatorTok{);}

    \CommentTok{// spawn α and β at gen + 1}
    \DataTypeTok{int}\NormalTok{ alpha }\OperatorTok{=}\NormalTok{ spawn\_child}\OperatorTok{(}\NormalTok{p\_snap}\OperatorTok{,}\NormalTok{ gen}\OperatorTok{=+}\DecValTok{1}\OperatorTok{);}
    \DataTypeTok{int}\NormalTok{ beta  }\OperatorTok{=}\NormalTok{ spawn\_child}\OperatorTok{(}\NormalTok{p\_snap}\OperatorTok{,}\NormalTok{ gen}\OperatorTok{=+}\DecValTok{1}\OperatorTok{);}

    \CommentTok{// reroute incident tunnels by polarity sign}
    \ControlFlowTok{for} \OperatorTok{(}\KeywordTok{auto}\OperatorTok{\&}\NormalTok{ t }\OperatorTok{:}\NormalTok{ tunnels}\OperatorTok{)} \OperatorTok{\{}
        \ControlFlowTok{if} \OperatorTok{(}\NormalTok{t}\OperatorTok{.}\NormalTok{cavity\_a }\OperatorTok{==}\NormalTok{ parent}\OperatorTok{)}\NormalTok{ t}\OperatorTok{.}\NormalTok{cavity\_a }\OperatorTok{=} \OperatorTok{(}\NormalTok{t}\OperatorTok{.}\NormalTok{polarity }\OperatorTok{\textgreater{}=} \DecValTok{0}\OperatorTok{)} \OperatorTok{?}\NormalTok{ alpha }\OperatorTok{:}\NormalTok{ beta}\OperatorTok{;}
        \ControlFlowTok{if} \OperatorTok{(}\NormalTok{t}\OperatorTok{.}\NormalTok{cavity\_b }\OperatorTok{==}\NormalTok{ parent}\OperatorTok{)}\NormalTok{ t}\OperatorTok{.}\NormalTok{cavity\_b }\OperatorTok{=} \OperatorTok{(}\NormalTok{t}\OperatorTok{.}\NormalTok{polarity }\OperatorTok{\textgreater{}=} \DecValTok{0}\OperatorTok{)} \OperatorTok{?}\NormalTok{ alpha }\OperatorTok{:}\NormalTok{ beta}\OperatorTok{;}
    \OperatorTok{\}}

    \CommentTok{// M4: binding tunnel}
    \ControlFlowTok{if} \OperatorTok{(}\NormalTok{rbg\_binding\_enabled}\OperatorTok{)}\NormalTok{ tunnels}\OperatorTok{.}\NormalTok{push\_back}\OperatorTok{(}\NormalTok{binding\_tunnel}\OperatorTok{(}\NormalTok{alpha}\OperatorTok{,}\NormalTok{ beta}\OperatorTok{));}

    \CommentTok{// M2: alias ledger update}
\NormalTok{    alias\_ledger}\OperatorTok{[}\NormalTok{parent}\OperatorTok{]} \OperatorTok{=} \OperatorTok{\{}\NormalTok{ RootAlias}\OperatorTok{,}\NormalTok{ parent}\OperatorTok{,} \OperatorTok{{-}}\DecValTok{1}\OperatorTok{,}\NormalTok{ parent\_label\_key }\OperatorTok{\};}
\NormalTok{    alias\_ledger}\OperatorTok{.}\NormalTok{push\_back}\OperatorTok{(\{}\NormalTok{ Alpha}\OperatorTok{,}\NormalTok{ parent}\OperatorTok{,}\NormalTok{ beta}\OperatorTok{,}\NormalTok{  parent\_label\_key }\OperatorTok{\});}
\NormalTok{    alias\_ledger}\OperatorTok{.}\NormalTok{push\_back}\OperatorTok{(\{}\NormalTok{ Beta}\OperatorTok{,}\NormalTok{  parent}\OperatorTok{,}\NormalTok{ alpha}\OperatorTok{,}\NormalTok{ parent\_label\_key }\OperatorTok{\});}

    \CommentTok{// cluster inheritance (children join parent\textquotesingle{}s HNCA cluster)}
    \OperatorTok{...}

\NormalTok{    fission\_map}\OperatorTok{[}\NormalTok{parent}\OperatorTok{]} \OperatorTok{=} \OperatorTok{\{}\NormalTok{alpha}\OperatorTok{,}\NormalTok{ beta}\OperatorTok{\};}

    \CommentTok{// M5: audit}
\NormalTok{    fission\_records}\OperatorTok{.}\NormalTok{push\_back}\OperatorTok{(}\NormalTok{record}\OperatorTok{(...));}
    \ControlFlowTok{return} \OperatorTok{\{}\NormalTok{alpha}\OperatorTok{,}\NormalTok{ beta}\OperatorTok{\};}
\OperatorTok{\}}
\end{Highlighting}
\end{Shaded}

Two variants are exposed: \texttt{fission()} acquires the physics mutex
and is safe to call from the shell; \texttt{fission\_locked()} assumes
the caller holds the mutex and is called from
\texttt{ConstructiveEngine::structural\_tick\_once()}.

\begin{center}\rule{0.5\linewidth}{0.5pt}\end{center}

\subsection{4. The Layer-2 Fission
Trigger}\label{the-layer-2-fission-trigger}

The trigger is a straight-line scalar loop, no branches inside
per-cavity computation:

\begin{Shaded}
\begin{Highlighting}[]
\ControlFlowTok{if} \OperatorTok{(}\NormalTok{fission\_enabled}\OperatorTok{)} \OperatorTok{\{}
    \DataTypeTok{int}\NormalTok{   fire\_id }\OperatorTok{=} \OperatorTok{{-}}\DecValTok{1}\OperatorTok{;}
    \DataTypeTok{float}\NormalTok{ fire\_I  }\OperatorTok{=} \FloatTok{0.0}\BuiltInTok{f}\OperatorTok{,}\NormalTok{ fire\_tau }\OperatorTok{=} \FloatTok{0.0}\BuiltInTok{f}\OperatorTok{;}
    \DataTypeTok{int}\NormalTok{   fire\_streak\_val }\OperatorTok{=}\NormalTok{ fission\_hold\_target }\OperatorTok{{-}} \DecValTok{1}\OperatorTok{;}
    \ControlFlowTok{for} \OperatorTok{(}\DataTypeTok{int}\NormalTok{ i }\OperatorTok{=} \DecValTok{0}\OperatorTok{;}\NormalTok{ i }\OperatorTok{\textless{}}\NormalTok{ N}\OperatorTok{;} \OperatorTok{++}\NormalTok{i}\OperatorTok{)} \OperatorTok{\{}
        \ControlFlowTok{if} \OperatorTok{(}\NormalTok{sub}\OperatorTok{{-}\textgreater{}}\NormalTok{cavities}\OperatorTok{[}\NormalTok{i}\OperatorTok{].}\NormalTok{generation }\OperatorTok{\textgreater{}=}\NormalTok{ sub}\OperatorTok{{-}\textgreater{}}\VariableTok{g\_max}\OperatorTok{)} \OperatorTok{\{}\NormalTok{ streak}\OperatorTok{[}\NormalTok{i}\OperatorTok{]=}\DecValTok{0}\OperatorTok{;} \ControlFlowTok{continue}\OperatorTok{;} \OperatorTok{\}}
        \ControlFlowTok{if} \OperatorTok{(}\NormalTok{sub}\OperatorTok{{-}\textgreater{}}\NormalTok{is\_fissioned}\OperatorTok{(}\NormalTok{i}\OperatorTok{))}                       \OperatorTok{\{}\NormalTok{ streak}\OperatorTok{[}\NormalTok{i}\OperatorTok{]=}\DecValTok{0}\OperatorTok{;} \ControlFlowTok{continue}\OperatorTok{;} \OperatorTok{\}}
        \DataTypeTok{float} \VariableTok{p\_} \OperatorTok{=}\NormalTok{ pos\_updates\_by\_cavity}\OperatorTok{[}\NormalTok{i}\OperatorTok{];}
        \DataTypeTok{float} \VariableTok{n\_} \OperatorTok{=}\NormalTok{ neg\_updates\_by\_cavity}\OperatorTok{[}\NormalTok{i}\OperatorTok{];}
        \DataTypeTok{float}\NormalTok{ denom }\OperatorTok{=} \VariableTok{p\_} \OperatorTok{+} \VariableTok{n\_} \OperatorTok{+} \FloatTok{1e{-}9}\BuiltInTok{f}\OperatorTok{;}
        \DataTypeTok{float}\NormalTok{ I\_i   }\OperatorTok{=} \OperatorTok{(}\VariableTok{p\_} \OperatorTok{*} \VariableTok{n\_}\OperatorTok{)} \OperatorTok{/} \OperatorTok{(}\NormalTok{denom }\OperatorTok{*}\NormalTok{ denom}\OperatorTok{);}          \CommentTok{// incoherence}
        \DataTypeTok{float}\NormalTok{ tau\_i }\OperatorTok{=}\NormalTok{ local\_tension}\OperatorTok{[}\NormalTok{i}\OperatorTok{];}
        \DataTypeTok{bool}\NormalTok{  trig  }\OperatorTok{=} \OperatorTok{(}\NormalTok{I\_i }\OperatorTok{\textgreater{}}\NormalTok{ fission\_I\_thresh}\OperatorTok{)} \OperatorTok{\&\&} \OperatorTok{(}\NormalTok{tau\_i }\OperatorTok{\textgreater{}}\NormalTok{ fission\_tau\_thresh}\OperatorTok{);}
\NormalTok{        streak}\OperatorTok{[}\NormalTok{i}\OperatorTok{]} \OperatorTok{=}\NormalTok{ trig }\OperatorTok{?}\NormalTok{ streak}\OperatorTok{[}\NormalTok{i}\OperatorTok{]} \OperatorTok{+} \DecValTok{1} \OperatorTok{:} \DecValTok{0}\OperatorTok{;}
        \ControlFlowTok{if} \OperatorTok{(}\NormalTok{streak}\OperatorTok{[}\NormalTok{i}\OperatorTok{]} \OperatorTok{\textgreater{}=}\NormalTok{ fission\_hold\_target }\OperatorTok{\&\&}\NormalTok{ streak}\OperatorTok{[}\NormalTok{i}\OperatorTok{]} \OperatorTok{\textgreater{}}\NormalTok{ fire\_streak\_val}\OperatorTok{)} \OperatorTok{\{}
\NormalTok{            fire\_streak\_val }\OperatorTok{=}\NormalTok{ streak}\OperatorTok{[}\NormalTok{i}\OperatorTok{];}
\NormalTok{            fire\_id }\OperatorTok{=}\NormalTok{ i}\OperatorTok{;}\NormalTok{ fire\_I }\OperatorTok{=}\NormalTok{ I\_i}\OperatorTok{;}\NormalTok{ fire\_tau }\OperatorTok{=}\NormalTok{ tau\_i}\OperatorTok{;}
        \OperatorTok{\}}
    \OperatorTok{\}}
    \ControlFlowTok{if} \OperatorTok{(}\NormalTok{fire\_id }\OperatorTok{\textgreater{}=} \DecValTok{0}\OperatorTok{)} \OperatorTok{\{}
\NormalTok{        sub}\OperatorTok{{-}\textgreater{}}\NormalTok{fission\_locked}\OperatorTok{(}\NormalTok{fire\_id}\OperatorTok{,}\NormalTok{ fire\_I}\OperatorTok{,}\NormalTok{ fire\_tau}\OperatorTok{,}\NormalTok{ fire\_streak\_val}\OperatorTok{,}\NormalTok{ tick}\OperatorTok{);}
        \CommentTok{// extend telemetry vectors for the two new cavities and one new tunnel}
        \OperatorTok{...}
    \OperatorTok{\}}
\OperatorTok{\}}
\end{Highlighting}
\end{Shaded}

\textbf{Design choices:}

\begin{itemize}
\tightlist
\item
  \textbf{Incoherence, not signed net.}
  \(I = \Sigma^+ \Sigma^- / (\Sigma^+ + \Sigma^-)^2\) peaks at 0.25 when
  positive and negative updates are exactly balanced, zero when one
  direction dominates. This is the destructive-interference signature RB
  targets.
\item
  \textbf{Local tension, not global.} Fission fires on cavities whose
  \emph{own} tension is high, not on cavities pulled up by a
  globally-tense IC. Prevents spurious fissions during high global
  stress.
\item
  \textbf{Streak requirement.} Three consecutive structural ticks with
  the trigger active. Blocks transient noise from firing.
\item
  \textbf{One fire per tick.} Even if multiple cavities cross the
  threshold, only the highest streak fires this tick. Prevents a single
  training sample from cascading through the graph.
\item
  \textbf{Generation gate.} Guarantees \(g_\text{max} = 1\) prevents
  grandchildren.
\end{itemize}

Aggregation vectors (\texttt{plasticity\_by\_cavity},
\texttt{pos\_updates\_by\_cavity}, \texttt{neg\_updates\_by\_cavity},
\texttt{corr\_ema}, \texttt{degree}, \texttt{local\_tension},
\texttt{cluster\_mean\_x}, \texttt{cluster\_mean\_pd}) are extended
in-place after fission fires. Old telemetry values are preserved; new
entries are zero-initialized; degree is recomputed by counting incident
tunnels.

\begin{center}\rule{0.5\linewidth}{0.5pt}\end{center}

\subsection{5. L6-A: Analogical Resonance
Bridge}\label{l6-a-analogical-resonance-bridge}

\subsubsection{5.1 Motivation}\label{motivation}

Fission on its own resolves \emph{within-concept} contradiction --- the
α child holds mechanistic-activation content, the β child holds
mechanistic-inhibition content. It does not create \emph{across-concept}
structure. In \texttt{l9.md} §3, L6 Analogical Mapping is the first
layer of true novelty: recognising that two disparate concepts share a
subgraph structure and forging a bridge that lets one carry information
about the other.

RB+G makes L6 physically tractable. After two related concepts fission,
their α children collect similar excitatory downstream targets
(metformin.α → ampk, berberine.α → ampk, resveratrol.α → ampk). The
analogy exists in the \emph{target overlap}. L6-A detects the overlap
and inserts a \emph{rubber tunnel} between the α children.

\subsubsection{5.2 Rubber Material and MIMS Impedance
Filtering}\label{rubber-material-and-mims-impedance-filtering}

The bridge material is drawn from \texttt{mims.md}: low stiffness
\(k = 0.05\) and heavy damping \(c = 0.15\) --- physical rubber. In MIMS
the material determines the frequency response of the tunnel.
High-frequency probes (20 Hz mechanistic pulses) get absorbed by the
low-stiffness edge (the small target cavity's inertia resists rapid
displacement, wave energy dissipates in damping). Low-frequency probes
(3 Hz structural pulses) transmit through --- a slow driver can move a
stiff or soft target equally well.

This gives the analogy tunnel a specificity property no symbolic
annotation can match: it is \emph{physically} a channel for structural
queries, physically silent for mechanistic queries. A 20 Hz probe on
\texttt{metformin} follows the direct \texttt{activates} and
\texttt{inhibits} tunnels (steel material, \(k \geq 1.6\)); the same
probe barely perturbs \texttt{berberine} via the rubber bridge. A 3 Hz
probe rides the rubber bridge freely.

\subsubsection{5.3 Discovery Algorithm}\label{discovery-algorithm}

\begin{Shaded}
\begin{Highlighting}[]
\DataTypeTok{int}\NormalTok{ discover\_analogies}\OperatorTok{(}\DataTypeTok{int}\NormalTok{ min\_shared }\OperatorTok{=} \DecValTok{2}\OperatorTok{,}
                       \DataTypeTok{float}\NormalTok{ rubber\_k }\OperatorTok{=} \FloatTok{0.05}\BuiltInTok{f}\OperatorTok{,} \DataTypeTok{float}\NormalTok{ rubber\_c }\OperatorTok{=} \FloatTok{0.15}\BuiltInTok{f}\OperatorTok{)}
\OperatorTok{\{}
\NormalTok{    analogy\_records}\OperatorTok{.}\NormalTok{clear}\OperatorTok{();}
    \BuiltInTok{std::}\NormalTok{vector}\OperatorTok{\textless{}}\DataTypeTok{int}\OperatorTok{\textgreater{}}\NormalTok{ parents }\OperatorTok{=}\NormalTok{ fissioned\_parents}\OperatorTok{();}     \CommentTok{// keys of fission\_map}
    \BuiltInTok{std::}\NormalTok{vector}\OperatorTok{\textless{}}\BuiltInTok{std::}\NormalTok{unordered\_set}\OperatorTok{\textless{}}\DataTypeTok{int}\OperatorTok{\textgreater{}\textgreater{}}\NormalTok{ targets}\OperatorTok{(}\NormalTok{parents}\OperatorTok{.}\NormalTok{size}\OperatorTok{());}
    \ControlFlowTok{for} \OperatorTok{(}\DataTypeTok{size\_t}\NormalTok{ i }\OperatorTok{=} \DecValTok{0}\OperatorTok{;}\NormalTok{ i }\OperatorTok{\textless{}}\NormalTok{ parents}\OperatorTok{.}\NormalTok{size}\OperatorTok{();} \OperatorTok{++}\NormalTok{i}\OperatorTok{)} \OperatorTok{\{}
        \DataTypeTok{int}\NormalTok{ a\_i }\OperatorTok{=}\NormalTok{ fission\_map}\OperatorTok{[}\NormalTok{parents}\OperatorTok{[}\NormalTok{i}\OperatorTok{]].}\NormalTok{first}\OperatorTok{;}        \CommentTok{// alpha id}
        \ControlFlowTok{for} \OperatorTok{(}\AttributeTok{const} \KeywordTok{auto}\OperatorTok{\&}\NormalTok{ t }\OperatorTok{:}\NormalTok{ tunnels}\OperatorTok{)} \OperatorTok{\{}
            \ControlFlowTok{if} \OperatorTok{(}\NormalTok{t}\OperatorTok{.}\NormalTok{polarity }\OperatorTok{\textless{}} \DecValTok{0}\OperatorTok{)} \ControlFlowTok{continue}\OperatorTok{;}               \CommentTok{// excitatory only}
            \ControlFlowTok{if} \OperatorTok{(}\NormalTok{t}\OperatorTok{.}\NormalTok{cavity\_a }\OperatorTok{==}\NormalTok{ a\_i}\OperatorTok{)}\NormalTok{      targets}\OperatorTok{[}\NormalTok{i}\OperatorTok{].}\NormalTok{insert}\OperatorTok{(}\NormalTok{t}\OperatorTok{.}\NormalTok{cavity\_b}\OperatorTok{);}
            \ControlFlowTok{else} \ControlFlowTok{if} \OperatorTok{(}\NormalTok{t}\OperatorTok{.}\NormalTok{cavity\_b }\OperatorTok{==}\NormalTok{ a\_i}\OperatorTok{)}\NormalTok{ targets}\OperatorTok{[}\NormalTok{i}\OperatorTok{].}\NormalTok{insert}\OperatorTok{(}\NormalTok{t}\OperatorTok{.}\NormalTok{cavity\_a}\OperatorTok{);}
        \OperatorTok{\}}
    \OperatorTok{\}}
    \ControlFlowTok{for} \OperatorTok{(}\DataTypeTok{size\_t}\NormalTok{ i }\OperatorTok{=} \DecValTok{0}\OperatorTok{;}\NormalTok{ i }\OperatorTok{\textless{}}\NormalTok{ parents}\OperatorTok{.}\NormalTok{size}\OperatorTok{();} \OperatorTok{++}\NormalTok{i}\OperatorTok{)}
      \ControlFlowTok{for} \OperatorTok{(}\DataTypeTok{size\_t}\NormalTok{ j }\OperatorTok{=}\NormalTok{ i}\OperatorTok{+}\DecValTok{1}\OperatorTok{;}\NormalTok{ j }\OperatorTok{\textless{}}\NormalTok{ parents}\OperatorTok{.}\NormalTok{size}\OperatorTok{();} \OperatorTok{++}\NormalTok{j}\OperatorTok{)} \OperatorTok{\{}
        \BuiltInTok{std::}\NormalTok{vector}\OperatorTok{\textless{}}\DataTypeTok{int}\OperatorTok{\textgreater{}}\NormalTok{ shared}\OperatorTok{;}
        \ControlFlowTok{for} \OperatorTok{(}\DataTypeTok{int}\NormalTok{ x }\OperatorTok{:}\NormalTok{ targets}\OperatorTok{[}\NormalTok{i}\OperatorTok{])} \ControlFlowTok{if} \OperatorTok{(}\NormalTok{targets}\OperatorTok{[}\NormalTok{j}\OperatorTok{].}\NormalTok{count}\OperatorTok{(}\NormalTok{x}\OperatorTok{))}\NormalTok{ shared}\OperatorTok{.}\NormalTok{push\_back}\OperatorTok{(}\NormalTok{x}\OperatorTok{);}
        \ControlFlowTok{if} \OperatorTok{((}\DataTypeTok{int}\OperatorTok{)}\NormalTok{shared}\OperatorTok{.}\NormalTok{size}\OperatorTok{()} \OperatorTok{\textless{}}\NormalTok{ min\_shared}\OperatorTok{)} \ControlFlowTok{continue}\OperatorTok{;}
\NormalTok{        insert\_rubber\_tunnel}\OperatorTok{(}\NormalTok{fission\_map}\OperatorTok{[}\NormalTok{parents}\OperatorTok{[}\NormalTok{i}\OperatorTok{]].}\NormalTok{first}\OperatorTok{,}
\NormalTok{                             fission\_map}\OperatorTok{[}\NormalTok{parents}\OperatorTok{[}\NormalTok{j}\OperatorTok{]].}\NormalTok{first}\OperatorTok{,}
\NormalTok{                             rubber\_k}\OperatorTok{,}\NormalTok{ rubber\_c}\OperatorTok{);}
\NormalTok{        analogy\_records}\OperatorTok{.}\NormalTok{push\_back}\OperatorTok{(\{...\});}
    \OperatorTok{\}}
    \ControlFlowTok{return}\NormalTok{ analogy\_records}\OperatorTok{.}\NormalTok{size}\OperatorTok{();}
\OperatorTok{\}}
\end{Highlighting}
\end{Shaded}

The intersection is over \emph{excitatory} targets only ---
inhibitory-target overlap does not produce a meaningful analogy under
this framework. \texttt{min\_shared\ =\ 2} is a low threshold; the demo
of §6.3 shows this admits some noise
(\texttt{molecule\ \textasciitilde{}\ ampk} with empty target list
appears because both cavities have external excitatory connections we
did not intersect). We report noise honestly rather than raising the
threshold to hide it.

\subsubsection{5.4 Repurpose Query}\label{repurpose-query}

\begin{Shaded}
\begin{Highlighting}[]
\DataTypeTok{void}\NormalTok{ repurpose}\OperatorTok{(}\DataTypeTok{int}\NormalTok{ label\_id}\OperatorTok{,} \DataTypeTok{float}\NormalTok{ low\_freq }\OperatorTok{=} \FloatTok{3.0}\BuiltInTok{f}\OperatorTok{,} \DataTypeTok{int}\NormalTok{ top\_k }\OperatorTok{=} \DecValTok{10}\OperatorTok{)} \OperatorTok{\{}
\NormalTok{    physics}\OperatorTok{.}\NormalTok{reset\_field}\OperatorTok{();}
\NormalTok{    physics}\OperatorTok{.}\NormalTok{drive}\OperatorTok{(}\NormalTok{label\_id}\OperatorTok{,}\NormalTok{ low\_freq}\OperatorTok{,} \FloatTok{6.0}\BuiltInTok{f}\OperatorTok{,} \DecValTok{20}\OperatorTok{);}
\NormalTok{    physics}\OperatorTok{.}\NormalTok{run}\OperatorTok{(}\DecValTok{120}\OperatorTok{);}
    \CommentTok{// rank all cavities by |displacement|, excluding the target itself}
    \CommentTok{// and its α/β children (they hold the input signal, not the analogy)}
    \KeywordTok{auto}\NormalTok{ ranked }\OperatorTok{=}\NormalTok{ sort\_by\_amplitude}\OperatorTok{(}\NormalTok{physics}\OperatorTok{.}\NormalTok{cavities}\OperatorTok{);}
\NormalTok{    print\_top}\OperatorTok{(}\NormalTok{ranked}\OperatorTok{,}\NormalTok{ top\_k}\OperatorTok{);}
\OperatorTok{\}}
\end{Highlighting}
\end{Shaded}

The output is a ranked list of cavities whose displacement is highest
under low-frequency drive on the target. Analogical partners (concepts
sharing structural role) rank ahead of mechanistic downstream targets
(which respond weakly at 3 Hz because their steel-material tunnels do
not transmit low frequencies well).

\begin{center}\rule{0.5\linewidth}{0.5pt}\end{center}

\subsection{6. Empirical Program}\label{empirical-program}

\subsubsection{6.1 20-Seed Reproducibility}\label{seed-reproducibility}

\textbf{Protocol.} For each label \(L\) in \{metformin, berberine,
atorvastatin, luna, linux, windows\_os\}:

\begin{itemize}
\tightlist
\item
  20 seeds drawn from \texttt{0xCAF17E00\ \^{}\ (0x9E3779B9\ *\ (i+1))}
  for \(i \in [0, 20)\).
\item
  Per seed, two passes: \textbf{BASELINE} with
  \texttt{fission\_enabled\ =\ false}; \textbf{RB+G} with all five
  mechanisms of §3.2 active.
\item
  Training: rounds = 5, NT profile \((D, C, S) = (1, 1, 2)\). Each round
  drives \(L\) at 20 Hz amplitude 8 for 40 ticks then runs 10 chunks of
  (20 physics ticks + 1 structural tick).
\item
  Probe: after restoring \(S = 1\), \texttt{reset\_field}, drive \(L\)
  at 20 Hz amplitude 6 for 15 ticks, run 60 ticks, measure global peak
  displacement.
\item
  Report per-seed rows and mean ± std across seeds.
\end{itemize}

\textbf{Aggregate results.}

\begin{longtable}[]{@{}
  >{\raggedright\arraybackslash}p{(\linewidth - 8\tabcolsep) * \real{0.2000}}
  >{\raggedright\arraybackslash}p{(\linewidth - 8\tabcolsep) * \real{0.2000}}
  >{\raggedright\arraybackslash}p{(\linewidth - 8\tabcolsep) * \real{0.2000}}
  >{\raggedright\arraybackslash}p{(\linewidth - 8\tabcolsep) * \real{0.2000}}
  >{\raggedright\arraybackslash}p{(\linewidth - 8\tabcolsep) * \real{0.2000}}@{}}
\toprule\noalign{}
\begin{minipage}[b]{\linewidth}\raggedright
Concept
\end{minipage} & \begin{minipage}[b]{\linewidth}\raggedright
Baseline gain
\end{minipage} & \begin{minipage}[b]{\linewidth}\raggedright
RB+G gain (RMS)
\end{minipage} & \begin{minipage}[b]{\linewidth}\raggedright
Fissions/seed
\end{minipage} & \begin{minipage}[b]{\linewidth}\raggedright
Target fissioned
\end{minipage} \\
\midrule\noalign{}
\endhead
\bottomrule\noalign{}
\endlastfoot
metformin (dual-tension) & 0.887 ± 0.045 & 1.517 ± 0.201 & 8.85 ± 0.96 &
20/20 \\
berberine (tri-hub) & 0.900 ± 0.030 & 1.515 ± 0.139 & 8.60 ± 1.16 &
20/20 \\
atorvastatin (mono-directional) & 1.007 ± 0.004 & 0.957 ± 0.034 & 9.75 ±
1.84 & 20/20 \\
luna (mutual + inhibition) & 1.244 ± 0.064 & 1.865 ± 0.162 & 7.00 ± 0.95
& 20/20 \\
linux (null: is\_a, has) & 1.032 ± 0.007 & 1.039 ± 0.057 & 5.00 ± 0.00 &
20/20 \\
windows\_os (pure null: is\_a) & 0.969 ± 0.008 & 0.955 ± 0.010 & 2.00 ±
0.00 & 0/20 \\
\end{longtable}

\textbf{Reading.}

\begin{itemize}
\tightlist
\item
  \textbf{Contradiction-cluster concepts} (metformin, berberine, luna)
  show large positive RB+G effect (Δ = +0.63 to +1.11 in RMS gain).
  Fission fires 7--9 times per seed; the target concept is fissioned in
  all 20 seeds.
\item
  \textbf{Mono-directional control} (atorvastatin) shows an \emph{honest
  null-adjacent effect}. Fission fires \textasciitilde10 times per seed
  (many on other high-tension nodes), but atorvastatin's own recall does
  not improve (RMS 0.957 vs 1.007 baseline --- actually marginally
  worse). Interpretation: fission on a concept without polarity conflict
  splits energy without accuracy benefit. RB+G is not a free-lunch
  amplifier.
\item
  \textbf{Null-adjacent control} (linux) shows a small positive global
  effect (1.38× vs 1.03× baseline) but the concept-scoped RMS is flat
  (1.039). Fissions fire on 5 unrelated nodes per seed; linux itself is
  fissioned in 20/20 (because its \texttt{is\_a\ operating\_system} +
  \texttt{is\_a\ kernel} edges have opposite MIMS materials, producing a
  small but nonzero incoherence).
\item
  \textbf{Pure null} (windows\_os) shows the cleanest control. The
  target is \emph{never} fissioned (0/20). Fissions fire on 2 unrelated
  nodes per seed (\texttt{operating\_system}, \texttt{software} --- hub
  concepts). The RB+G gain is 0.955 vs 0.969 baseline --- statistically
  indistinguishable.
\end{itemize}

The mechanism \emph{refuses to fire} when the substrate has no
contradiction to resolve. This is the falsifiability signature.

\subsubsection{6.2 Four-Lever Ablation
Table}\label{four-lever-ablation-table}

\textbf{Protocol.} For metformin and luna, 5 seeds each, six
configurations toggling one RB+G mechanism at a time. Reported: concept
RMS gain (post/pre using target and α/β displacement) and fissions per
seed.

\textbf{Metformin (n=5).}

\begin{longtable}[]{@{}
  >{\raggedright\arraybackslash}p{(\linewidth - 8\tabcolsep) * \real{0.2000}}
  >{\raggedright\arraybackslash}p{(\linewidth - 8\tabcolsep) * \real{0.2000}}
  >{\raggedright\arraybackslash}p{(\linewidth - 8\tabcolsep) * \real{0.2000}}
  >{\raggedright\arraybackslash}p{(\linewidth - 8\tabcolsep) * \real{0.2000}}
  >{\raggedright\arraybackslash}p{(\linewidth - 8\tabcolsep) * \real{0.2000}}@{}}
\toprule\noalign{}
\begin{minipage}[b]{\linewidth}\raggedright
Configuration
\end{minipage} & \begin{minipage}[b]{\linewidth}\raggedright
RMS gain
\end{minipage} & \begin{minipage}[b]{\linewidth}\raggedright
Fissions
\end{minipage} & \begin{minipage}[b]{\linewidth}\raggedright
Target fissioned
\end{minipage} & \begin{minipage}[b]{\linewidth}\raggedright
Reading
\end{minipage} \\
\midrule\noalign{}
\endhead
\bottomrule\noalign{}
\endlastfoot
BASELINE (all off) & 0.922 ± 0.011 & 0.0 & 0/5 & negative control \\
ABLATE\_CAP (\(g_\text{max}=999\)) & 1.536 ± 0.239 & 25.6 ± 3.3 & 5/5 &
cascade; 3× fissions, no accuracy win \\
ABLATE\_ALIAS & 2.495 ± 0.144 & 3.4 ± 0.8 & 5/5 & probe hits ghost
cavity --- measurement confound \\
ABLATE\_FANOUT (factor = 1.0) & 1.856 ± 0.069 & 13.6 ± 1.4 & 5/5 &
\(\sqrt{2}\times\) energy injection --- inflated peak \\
ABLATE\_BIND (no binding tunnel) & 1.470 ± 0.057 & 8.8 ± 0.8 & 5/5 &
minor loss vs FULL \\
FULL (all five mechanisms) & 1.345 ± 0.054 & 8.4 ± 1.0 & 5/5 & RB+G \\
\end{longtable}

\textbf{Luna (n=5).}

\begin{longtable}[]{@{}llll@{}}
\toprule\noalign{}
Configuration & RMS gain & Fissions & Target fissioned \\
\midrule\noalign{}
\endhead
\bottomrule\noalign{}
\endlastfoot
BASELINE & 0.489 ± 0.124 & 0.0 & 0/5 \\
ABLATE\_CAP & 3.203 ± 0.319 & 30.2 ± 3.3 & 5/5 \\
ABLATE\_ALIAS & 3.637 ± 0.394 & 4.6 ± 0.8 & 5/5 \\
ABLATE\_FANOUT & 2.987 ± 0.324 & 7.4 ± 0.5 & 5/5 \\
ABLATE\_BIND & 1.904 ± 0.134 & 6.8 ± 0.8 & 5/5 \\
FULL & 1.914 ± 0.152 & 7.0 ± 0.9 & 5/5 \\
\end{longtable}

\textbf{Honest reading.} Two ablations (\texttt{ABLATE\_ALIAS},
\texttt{ABLATE\_FANOUT}) produce \emph{higher} raw RMS than FULL. We do
not hide this.

\begin{itemize}
\tightlist
\item
  \textbf{ABLATE\_ALIAS} disables \texttt{rbg\_alias\_route\_enabled}.
  The probe on the fissioned parent hits the zeroed ghost cavity and
  never reaches the children. The recorded pre and post peaks measure
  neighbour drift, not target recall. The high number reflects a
  \emph{change in measurement semantics}, not a real accuracy gain. This
  is a \textbf{documented confound}; the ablation cannot be cleanly
  interpreted, which is precisely why alias routing is a mechanism.
\item
  \textbf{ABLATE\_FANOUT} sets the fan-out factor to 1.0 instead of
  \(1/\sqrt{2}\). The two children each receive full drive amplitude;
  total kinetic energy injected is \(2\times\) a single-cavity probe.
  The RMS peak inflates by \(\sqrt{2}\) purely from extra energy, not
  from better learning. This is an \textbf{honest-accounting failure},
  and the reason the fan-out is required.
\item
  \textbf{ABLATE\_CAP} allows cascade. Fissions rise 3--4× without a
  corresponding gain --- the mechanism becomes computationally expensive
  without improving concept recall. Grandchildren of grandchildren carry
  vanishingly small amplitude; the concept becomes hard to interpret.
\item
  \textbf{ABLATE\_BIND} removes the α↔β spring. Loss vs FULL is small
  (metformin: 1.47 vs 1.35 --- actually a marginal \emph{gain} in this
  seed set; luna: 1.90 vs 1.91 --- flat). Evidence that RB's fiat
  \(k=2.5\) was over-specified: a much softer learned coupling (or none
  at all) works nearly as well. We retain the binding tunnel as a
  coordination default; a version without it is a valid variant.
\end{itemize}

\textbf{What FULL buys} is not the biggest number. It buys: (i) honest
energy accounting, (ii) cascade-free convergence, (iii) label-clean
audit trails. The number is the number; the mechanism protects the
interpretation.

\subsubsection{6.3 L6 Analogical Bridging --- Drug Repurposing Case
Study}\label{l6-analogical-bridging-drug-repurposing-case-study}

\textbf{Setup.} Insulin-signalling ontology from
\texttt{data/insulin\_pathway.axioms}. Train \texttt{metformin} then
\texttt{berberine} via \texttt{train\_multi\ metformin,berberine\ 5}.
Run \texttt{discover\_analogies\ 2}. Then
\texttt{repurpose\ metformin\ 3.0\ 15}.

\textbf{Discovered analogies} (five bridges inserted):

\begin{longtable}[]{@{}
  >{\raggedright\arraybackslash}p{(\linewidth - 2\tabcolsep) * \real{0.5000}}
  >{\raggedright\arraybackslash}p{(\linewidth - 2\tabcolsep) * \real{0.5000}}@{}}
\toprule\noalign{}
\begin{minipage}[b]{\linewidth}\raggedright
Analogy
\end{minipage} & \begin{minipage}[b]{\linewidth}\raggedright
Shared excitatory targets
\end{minipage} \\
\midrule\noalign{}
\endhead
\bottomrule\noalign{}
\endlastfoot
\texttt{akt\ \textasciitilde{}\ ampk} & glut4, kinase \\
\texttt{akt\ \textasciitilde{}\ mtor} & insulin\_signaling, kinase \\
\texttt{type2\_diabetes\ \textasciitilde{}\ drug} & pioglitazone,
liraglutide, sitagliptin, canagliflozin, empagliflozin (6 total) \\
\texttt{berberine\ \textasciitilde{}\ resveratrol} & (via base graph) \\
\texttt{molecule\ \textasciitilde{}\ ampk} & (empty target list ---
reported as noise) \\
\end{longtable}

Real analogies dominate:

\begin{itemize}
\tightlist
\item
  \textbf{\texttt{akt\ \textasciitilde{}\ ampk}} --- both are kinases
  activating \texttt{glut4}; both play insulin-cascade roles. Real
  biology.
\item
  \textbf{\texttt{akt\ \textasciitilde{}\ mtor}} --- sequential kinases
  in the insulin-signalling pathway. Real biology.
\item
  \textbf{\texttt{berberine\ \textasciitilde{}\ resveratrol}} --- both
  natural-product AMPK activators. \textbf{A known clinical
  drug-repurposing pair}.
\item
  \textbf{\texttt{type2\_diabetes\ \textasciitilde{}\ drug}} ---
  cross-category structural analogy, six shared drug targets.
\end{itemize}

\textbf{\texttt{metformin} at 3 Hz --- top 15 vibrating cavities.}

\begin{longtable}[]{@{}lll@{}}
\toprule\noalign{}
Rank & Cavity & Displacement \\
\midrule\noalign{}
\endhead
\bottomrule\noalign{}
\endlastfoot
1 & mtor.β & −0.09787 \\
2 & rapamycin.β & +0.05202 \\
3 & ampk.α & +0.03914 \\
4 & berberine.α & +0.03419 \\
5 & hmg\_coa\_reductase.β & −0.01533 \\
6 & energy\_homeostasis & +0.01509 \\
7 & resveratrol.α & +0.01270 \\
8 & berberine.β & +0.01087 \\
9 & obesity & +0.00970 \\
10 & \textbf{atorvastatin} & \textbf{+0.01486} \\
\ldots{} & \ldots{} & \ldots{} \\
\end{longtable}

\textbf{The key row is rank 10.} In a negative control run that fissions
both concepts but skips \texttt{discover\_analogies},
\texttt{atorvastatin} does \emph{not} appear in the top 15. In the naive
run (no fission, no bridges) \texttt{atorvastatin} appears at rank
\textasciitilde10 with amplitude 0.014 --- but only because of the
direct edge \texttt{atorvastatin\ inhibits\ hmg\_coa\_reductase} which
is a 2-hop path from berberine to atorvastatin.

The L6-A signal is: after fission + bridge insertion, the wave path
\texttt{metformin\ →\ metformin.α\ →\ (rubber\ bridge)\ →\ berberine.α\ →\ hmg\_coa\_reductase\ →\ atorvastatin}
carries measurable amplitude that reaches atorvastatin. No axiom says
\texttt{metformin\ \textasciitilde{}\ berberine}; the rubber bridge did
not exist until \texttt{discover\_analogies} created it.

\subsubsection{6.4 Storage-Floor Honesty}\label{storage-floor-honesty}

An honest cognitive architecture names its storage. By the Translation
Capacity Lemma {[}Mahir 2026p, §3.1{]} no system that reliably
discriminates \(N\) discrete answers escapes storing \(\log_2 N\) bits.

\textbf{What is stored} (primitive facts,
\texttt{data/insulin\_pathway.axioms}, \textasciitilde50 lines):

\begin{itemize}
\tightlist
\item
  \texttt{metformin\ activates\ ampk}
\item
  \texttt{berberine\ activates\ ampk}
\item
  \texttt{atorvastatin\ inhibits\ hmg\_coa\_reductase}
\item
  47 other atomic drug--target and pathway facts
\end{itemize}

\textbf{What is NOT stored anywhere in the pipeline:}

\begin{itemize}
\tightlist
\item
  The analogy \texttt{metformin\ \textasciitilde{}\ berberine} (no axiom
  line, no relation type in the ontology).
\item
  The child cavities \texttt{metformin.α}, \texttt{berberine.α} (spawned
  by physics at training time).
\item
  The rubber bridge tunnel (inserted by \texttt{discover\_analogies} at
  runtime).
\item
  The answer that \texttt{atorvastatin} is repurposable from
  \texttt{metformin} (constructed by wave propagation through the
  physics-inserted bridge).
\item
  No \texttt{(query,\ answer)} pair anywhere in the training data.
\end{itemize}

The 50 axiom lines are \textasciitilde1 KB. The reachable derivation
space under RB+G + L6-A grows exponentially with the number of fissioned
concepts. This is the composition-over-retrieval definition of §3.4 of
the SI paper --- a small primitive set producing an exponentially larger
reachable derivation set.

This is also how a human brain works. A child who never reads
\texttt{metformin\ is\ a\ drug} cannot dream about it. Storing
primitives is not the failure mode; storing \emph{answers} while
pretending to derive them is.

\begin{center}\rule{0.5\linewidth}{0.5pt}\end{center}

\subsection{7. Comparison to Prior Work}\label{comparison-to-prior-work}

\textbf{vs.~LLM next-token prediction.}

\begin{longtable}[]{@{}
  >{\raggedright\arraybackslash}p{(\linewidth - 4\tabcolsep) * \real{0.3333}}
  >{\raggedright\arraybackslash}p{(\linewidth - 4\tabcolsep) * \real{0.3333}}
  >{\raggedright\arraybackslash}p{(\linewidth - 4\tabcolsep) * \real{0.3333}}@{}}
\toprule\noalign{}
\begin{minipage}[b]{\linewidth}\raggedright
\end{minipage} & \begin{minipage}[b]{\linewidth}\raggedright
LLM
\end{minipage} & \begin{minipage}[b]{\linewidth}\raggedright
RB+G
\end{minipage} \\
\midrule\noalign{}
\endhead
\bottomrule\noalign{}
\endlastfoot
Semantics & Distributed across \(\sim 10^{12}\) weights & Standing-wave
field on \(\sim 200\) cavities \\
Novelty & Interpolation over memorised \((Q, A)\) pairs & Fission +
bridge insertion + wave propagation \\
Storage & Answer pairs baked into weights & Only primitive facts;
\textasciitilde50 axiom lines \\
Repurposing & Requires the pair to have been in the corpus & Emerges
from physics-computed target intersection \\
Failure mode & Hallucination & Fission refuses to fire (windows\_os:
0/20) \\
Explainability & None --- post-hoc attention weights & Fission ledger +
analogy record are the derivation \\
\end{longtable}

\textbf{vs.~original RB {[}Mahir 2026o{]}.} RB proposes fission at
\(I > 0.15\), \(\tau\) threshold, with binding at fiat \(k = 2.5\). Our
attempt hit cascade (F1), timeout (F2), null labels (F3), drive confound
(F4). RB+G fixes each: generation cap, alias ledger, \(1/\sqrt{2}\)
fan-out, learned binding, audit trail. Ablation shows the fiat binding
was over-specified --- a much softer coupling suffices.

\textbf{vs.~SRG {[}Mahir 2026t{]}.} SRG keeps the cavity as one entity
and routes forces into E/I ports. Empirically SRG produced marginal
regression on our metformin benchmark (0.887 → 0.831). SRG and RB+G are
compositional: SRG can run alongside RB+G (both mechanisms are guarded
by their own toggles in the physics engine). We do not use SRG in the
results reported here.

\textbf{vs.~CMA bisociation {[}Mahir 2026u{]}.} CMA forces edges between
semantically distant nodes and resolves the tension via intermediate
hypothesis paths. RB+G is orthogonal: it splits \emph{within} a concept
(self-contradiction), while CMA bridges \emph{across} concepts (external
contradiction). L6-A is the natural bridge between them --- after RB+G
fissions two concepts, L6-A inserts CMA-style structural bridges between
their α children.

\textbf{vs.~Layer-2 nine primitives {[}Mahir 2026p, §4.4{]}.} The nine
reasoning primitives (deduction, transitive composition, inheritance,
abduction, negation, case-analysis, induction, analogy, recursion) are
graph-rewriting rules on a fixed vertex set. RB+G is a rule that
\emph{grows} the vertex set --- a realization of Axiom II
(non-equilibrium phase-space expansion). It is not a tenth primitive; it
is a substrate operation the primitives cannot express, because they
cannot spawn new vertices.

\begin{center}\rule{0.5\linewidth}{0.5pt}\end{center}

\subsection{8. Limitations}\label{limitations}

\begin{itemize}
\tightlist
\item
  \textbf{Ontology scale.} All results are on ontologies of
  \textasciitilde200 concepts. Scaling to \(10^4\)--\(10^5\) concept
  graphs will require: (i) sparse aggregation in the structural tick
  (currently \(O(N + T)\) per tick), (ii) fission trigger evaluated only
  on cavities that touched a driven region in the recent past. Not
  tested.
\item
  \textbf{Fission-trigger tuning.} The thresholds \(I > 0.15\),
  \(\tau > 0.005\), streak \(= 3\) are set by inspection on the
  insulin-pathway ontology. A domain-independent auto-tuning rule
  (e.g.~based on quantiles of the cavity-wise \(I\) distribution) has
  not been demonstrated.
\item
  \textbf{L6-A discovery is post-hoc.} \texttt{discover\_analogies} runs
  after training. An online variant that inserts bridges continuously as
  α-children accumulate shared targets is a natural extension, not
  implemented here.
\item
  \textbf{Rubber-tunnel frequency is a design choice.} The claim that
  low-frequency probes ride the bridge while high-frequency probes are
  absorbed is a consequence of the material impedance mismatch between
  rubber (\(k=0.05\)) and steel (\(k \geq 1.6\)) tunnels. The optimal
  probe frequency is domain-dependent; we use 3 Hz throughout for the
  insulin ontology.
\item
  \textbf{Storage floor is not eliminated.} By the Translation Capacity
  Lemma this cannot be otherwise; we name it explicitly. What we claim
  is a \emph{low} storage floor (\textasciitilde50 axioms →
  exponentially larger reachable derivation set), not zero.
\item
  \textbf{The two ``high raw RMS'' ablations are honest confounds.} We
  report them as such; they do not falsify the mechanism.
\end{itemize}

\begin{center}\rule{0.5\linewidth}{0.5pt}\end{center}

\subsection{9. Conclusion}\label{conclusion}

RB+G shows that a physics substrate can \emph{grow its own state space}
under sustained contradiction, without any global gradient signal, and
without pattern-matching over stored \((Q, A)\) pairs. The mechanism is
disciplined: five specific fixes address four documented failure modes
of the underlying RB proposal, and every fission event is auditable via
the \texttt{data/fission\_ledger.jsonl} trail. On six concepts spanning
contradiction-cluster, mono-directional, and pure-null cases, the
mechanism produces large positive effects where it should and honest
null effects where it should not.

L6-A closes the loop from within-concept splitting to across-concept
structural analogy: after two related concepts fission, the intersection
of their α-children's downstream excitatory targets is used to insert a
rubber tunnel between the α cavities. On the insulin-signalling ontology
this surfaces real drug-repurposing pairs
(\texttt{berberine\ \textasciitilde{}\ resveratrol}) and reveals
cross-concept reachability (\texttt{metformin} → \texttt{atorvastatin}
via the physics-inserted bridge) that no path in the raw axiom graph
provides.

The storage floor is 50 axiom lines. The reachable derivation set
includes every fission spawn, every analogy bridge, and every
wave-propagation answer. This is the composition-over-retrieval
definition of a thinker --- not because the primitives are absent, but
because their compositions dominate.

\begin{center}\rule{0.5\linewidth}{0.5pt}\end{center}

\subsection{References}\label{references}

{[}Mahir 2026d{]} The ZeroGrad Cognitive Engine.
\texttt{research-paper/zerograd.md}. {[}Mahir 2026g{]} Recursive
Integrity Field Architecture (RIFA). \texttt{research-paper/rifa.md}.
{[}Mahir 2026m{]} Bounded Self-Model Architecture (BSMA).
\texttt{research-paper/bmsa\_new.md},
\texttt{research-paper/phenomenal-consciousness.md}. {[}Mahir 2026n{]}
Hierarchical NCA Constraint Solver (HNCA).
\texttt{research-paper/hnca.md}. {[}Mahir 2026o{]} Resonant Bifurcation
(Node Fission). \texttt{research-paper/rb.md}. {[}Mahir 2026p{]} From
Artificial to Synthetic Intelligence.
\texttt{research-paper/synthetic-intelligence.md}. {[}Mahir 2026q{]} The
Phonic Resonance Matrix (PRM). \texttt{research-paper/prm.md}. {[}Mahir
2026r{]} Mechanical Impedance and Material Stiffness (MIMS).
\texttt{research-paper/mims.md}. {[}Mahir 2026s{]} Neuro-Topological
Dynamics Architecture (NTDA). \texttt{research-paper/ntda.md}. {[}Mahir
2026t{]} Synaptic Receptor Gating (SRG). \texttt{research-paper/srg.md}.
{[}Mahir 2026u{]} Conceptual Mutation Architecture (CMA ---
Bisociation). \texttt{research-paper/CMA.md}. {[}Koestler 1964{]}
\emph{The Act of Creation}. Hutchinson. {[}Chomsky 1957{]} Syntactic
Structures. Mouton.

\begin{center}\rule{0.5\linewidth}{0.5pt}\end{center}

\textbf{RB+G: contradiction metabolism through disciplined bifurcation;
analogy through physics-inserted rubber bridges; construction over
retrieval, and storage where storage cannot be avoided.}

\end{document}
